
\[
 \ch(i_*(i^*(y)))=\ch(y)\cdot\ch(1-J).
\]
Comme \(\underline{J}\) est inversible, on a \(\widetilde c(J)=(1,1+c^1(J))\), donc \(\ch(J)=\exp(c^1(J))\).
Par suite, on a :
\[
 \ch(i_*(i^*(y)))=\ch(y)\cdot(1-\exp(c^1(J))).
\]

Par ailleurs, la fonctorialité de \(\ch\), et la formule de projection 3.7 entrainent :
\[
 i_{\mathrm{gr}}\bigl(\ch(i^*(y))\cdot\Todd(\check T_i)\bigr)
 =
 \ch(y)\cdot i_{\mathrm{gr}}(\Todd(\check T_i)).
\]
Or \(T_i=-\cl(\underline{J}/\underline{J}^2)=-i^*(J)\). Donc :
\[
 i_{\mathrm{gr}}(\Todd(\check T_i))
 =
 i_{\mathrm{gr}}\bigl(i^{\mathrm{gr}}(\Todd(-\check J))\bigr)
 =
 i_{\mathrm{gr}}(1)\cdot\Todd(-\check J).
\]

Comme \(i_*(1)=1-J\), \(i_{\mathrm{gr}}(1)\) est l'image dans \(\Gr^1K^{\bullet}(X)\) de \(1-J\), donc c'est \(-c(J)\). Par ailleurs, on a :
\[
 \Todd(-\check J)=(\Todd(\check J))^{-1}
 =
 \frac{1-\exp(c^1(J))}{-c^1(J)}.
\]

On obtient donc finalement :
\[
 i_{\mathrm{gr}}\bigl(\ch(i^*(y))\cdot\Todd(\check T_i)\bigr)
 =
 \ch(y)\cdot(1-\exp(c^1(J))),
\]
ce qui prouve la relation à démontrer.

\section*{5. Théorème de Riemann--Roch : cas du morphisme structural d'un fibré projectif}

\textbf{5.1.} D'après 3.9, il reste, pour prouver le théorème de Riemann--Roch, à le prouver pour le morphisme structural \(f:\underline{\underline{P}}^r_Y\longrightarrow Y\), où \(Y\) est un schéma quasi-compact, admettant un faisceau ample. On pose \(P=\underline{\underline{P}}^r_Y\), \(\xi=\cl(\underline{O}_P(1))\), et \(\underline{E}=(\underline{O}_Y)^{r+1}\).

